% Ad Familiares 9.16 (ad-familiares-09-16) --- English translation
% Translated by Claude (Anthropic) from the Latin (Perseus).
% Style guide: STYLE.md.

\ciceroLetterOpener{CICERO PAPIRIO PAETO S.}

\ciceroSection{1} Your letter gave me delight, and in it the first thing I loved was your love --- which had prompted you to write because you were afraid Silius's message might have brought me some worry. About that you had already written to me before, and twice indeed in identical terms, so that I could easily tell you were agitated, and I had written back with care, so as in such a matter and at such a time either to free you from that anxiety or at any rate to lighten it.

\ciceroSection{2} But since in your most recent letter as well you show how much the business weighs on you, hold this for certain, my Paetus: whatever could be done by art (for it is no longer enough to fight by strategy --- some species of artifice has to be devised), whatever could be worked out or carried through toward winning over and gathering in the goodwill of those people, I have pursued with the utmost zeal, and not, as I think, in vain; for I am so courted, so attended to, by all those who are loved by Caesar, that I suppose I am loved by them. Now true love and feigned love are not easily told apart, except when some moment of the right kind comes along, by which, as gold by fire, loyal goodwill can be tested by some peril. Other signs are signs in common, but I myself use one argument why I judge that I am loved truly and from the heart: it is that our position and theirs are such that they would have no reason to pretend.

\ciceroSection{3} As for the man in whose hand all power lies, I see nothing to fear from him --- except that all things are uncertain once law has been departed from, and nothing can be guaranteed as to what shape it will take, given that it rests on another's will, not to say his whim. Yet so far as he is concerned, his temper has been roused against me on no point; for in this very matter the greatest moderation has been used by me. As I once held that free speech belonged to me, by whose effort freedom existed in the state, so now that freedom is lost I hold that I should say nothing to offend either his will or that of those he loves. But if I wanted to escape the reputation of certain sharp or witty sayings, I should have to throw away my reputation for talent; and that, if I could do it, I would not refuse.

\ciceroSection{4} But for all that, Caesar himself has very acute judgement, and just as your brother Servius --- whom I count as the most lettered of men --- could readily say ``This verse is not Plautus's, this one is,'' because he had a worn ear for marking the classes of poets and from habitual reading, so I am told that Caesar, now that he has compiled whole volumes of sayings \textit{[Greek: apophthegmaton]}, makes a point of rejecting any that is brought to him as mine and is not. He does this all the more now, because his intimates live with me almost daily; and many things crop up in casual talk that perhaps strike them, when I have said them, as neither unlearned nor flat, and these are conveyed to him along with the rest of the day's business --- for that is what he himself directed. So it falls out that if he hears anything besides about me, he reckons it not worth hearing. For which reason I am making no use of your Oenomaus --- though you placed the lines of Accius well.

\ciceroSection{5} But what is this ``envy,'' and what is there now in me for anyone to envy? Yet grant it is all so. I see that this is the view of those philosophers who alone, as it seems to me, grasp the power of virtue: that the wise man has no warrant to give for anything but his own fault. From this I think I am twice clear --- both because I held the views that were most upright, and because, when I saw that there was not protection enough to hold them, I judged I should not contend in force with stronger men. So in the duty of the good citizen I am at any rate not to be blamed. What remains is that I should say nothing or do nothing rashly, nothing stupidly, against the powerful. That too I count as part of the wise man's part. As for the rest --- what anyone says I have said, or how he takes it, or with what loyalty those who assiduously court and attend on me live with me --- I cannot stand surety for that.

\ciceroSection{6} So it falls out that I console myself both by the consciousness of my earlier counsels and by the moderation of the present, and transfer that simile from Accius not just now to ``envy'' but to fortune --- which I judge a light and feeble thing, and which by a firm and grave mind ought to be broken ``like a wave against a rock.'' Indeed when the records of the Greeks are full of how the wisest men bore monarchy at Athens or at Syracuse --- being themselves in a manner of speaking free though their cities were enslaved --- shall I not think that I can hold my own ground in such a way as neither to offend any man's mind nor to break my own dignity?

\ciceroSection{7} Now I come to your jokes, since following the Oenomaus of Accius you brought on not, as it used to be, an Atellan farce, but, as the fashion now goes, a mime. What is this \textsuperscript{\textdagger}popillius\textsuperscript{\textdagger} you tell me of, what \textsuperscript{\textdagger}denarius\textsuperscript{\textdagger}, what dish of cheese-and-salt-fish? Such things were borne with by my easy nature once; now the case is altered. Hirtius and Dolabella I have as pupils in speaking, masters in dining; for I imagine you have heard --- if everything is reported up to your part of the country --- that they declaim at my house, I dine at theirs. Now as for your forswearing solvency to me, that is no good; for once, when you had means, your little ventures kept you on your toes, but now, when you are so equably parting with your goods, do not be of the mind that, when you take me in as a guest, you will think you are receiving some property valuation --- and even that is the lighter blow when it comes from a friend than from a debtor.

\ciceroSection{8} Nor, for all that, am I after the kind of dinner that leaves great leftovers; let what there is be sumptuous and elegant. I remember your telling me about Phamea's dinner. ``Let it be earlier,'' you said; for the rest, the same way. And if you persist in calling me back to your mother's dinner, I shall bear that too --- for I want to see the spirit that dares to set in front of me the things you describe, or indeed an octopus to match a red-leaded Jupiter. Believe me, you will not dare; before I arrive, the report of my new elegance will reach you, and you will quail at it. Nor is there any reason to pin some hope on the appetizer course; the whole of that I have done away with, for I used to be done in beforehand by your olives and your Lucanian sausages.

\ciceroSection{9} But why are we talking like this? Only let me get out there. As for you --- for I do want to wipe away your apprehension of mind --- go back to your good old cheese-and-salt-fish. The one expense I shall bring you is that you will have to heat the bath; for the rest, our usual way. The earlier talk was play.

\ciceroSection{10} About the Selician villa, you both saw to the matter diligently and wrote about it most amusingly; and so I think I shall pass it over. Salt there is enough; clowns, too few.
